\section{Eliminating interior area multipliers}
\label{var:singular-multipliers}

We prove that the area multiplier obtained from separation is concentrated on
the given boundary.  The main issue is that an area minimizer can have a
singular frontier in high dimension.  Atomic transverse mass forces a regular
sheet to be totally geodesic.  For the remaining sheets, a common positive
Jacobi field, null focusing, and a cone argument give a uniform curvature
bound.  A separate locality argument removes diffuse multiplier mass from
sheets on which the lapse vanishes.  All these arguments concern the original
smooth data.

Throughout this section, $k=n-1$, and $\mathcal F_*$ and $\pi$ are the compact
family and probability measure of Proposition~\ref{var:area-envelope} and
equation~\ref{var:separation}.  We use the full-measure family supplied by
Lemma~\ref{var:trace-free-sheets}, intersected with $\operatorname{supp}\pi$.
A member of this family has frontier $S$,
or has compact frontier wholly in the open exterior and satisfies
\begin{equation}\label{var:zero-sheet-expansions}
 H_\Sigma=0,\qquad \operatorname{tr}_\Sigma K=0
\end{equation}
on every free regular sheet.  We retain the local common-sheet,
noncrossing, and graphical compactness conclusions of
Proposition~\ref{var:area-envelope}.  In particular, distinct nearby sheets
are disjoint graphs, and on smaller patches the tilt between them is bounded
by a constant times their separation.

The moments of Proposition~\ref{var:moment-regularity} satisfy
\begin{equation}\label{var:singular-system}
 \operatorname{sym}\nabla X=-uK,\qquad
 \nabla^2u=B(u,X,\nabla X)+\mathcal M,
 \qquad
 \mathcal M=\frac c2\int_{\mathcal F_*}
       \nu_D^\flat\otimes\nu_D^\flat\,dA_D\,d\pi(D)
\end{equation}
in the open exterior.  Here $c>0$, $u$ is locally Lipschitz, $X$ is locally
$C^{1,1}$, and the last formula defines a tensor-valued measure relative to
the smooth volume density.  Its singular sets have zero sheet area.
Here $B$ abbreviates the operator $\mathcal L$ of
Proposition~\ref{var:moment-regularity}.  It is the precise lower-order
metric-adjoint expression in
equation~\ref{var:separation}, with the electric flux held fixed and the
unit-ball rays transported isometrically.  It has smooth coefficients and
is linear in its displayed arguments.  Its origin, and not merely its
regularity, will be used in the focusing calculation.  We also have
\begin{equation}\label{var:singular-complementarity}
 u\geq |X|,\qquad \mu_m\geq |J_m|,\qquad
 u\mu_m+J_m(X)=0.
\end{equation}

\subsection{Local leaf measures and atomic sheets}

Fix a free regular incidence point $(D,p)$.  The uniform graph conclusion
of Proposition~\ref{var:area-envelope}, after shrinking its neighborhood,
provides nested open spatial boxes
\[
 p\in W\Subset V\Subset U,
 \qquad W=B_0\times I_0,
\]
in coordinates $(y,t)$, and a larger base ball $B_1\supset\overline{B_0}$,
with the following property.  Every minimizing frontier that meets $V$
has in $U$ a unique selected graph $t=f_F(y)$ over $B_1$; all portions of
that frontier in $W$ belong to this graph.  Its graph over $B_1$ lies in
the coordinate chart, and the graphs have uniform interior smooth
bounds and are mutually disjoint unless they coincide.  The boxes can
be chosen with this property by first taking the uniform graph
neighborhood and then choosing the middle and inner boxes with positive
margin from its side and top faces.

Put
\[
 \mathcal H_V=\{F\in\mathcal F_*:
                         \operatorname{frontier}(F)\cap V\ne\varnothing\}.
\]
This is an open subset of the minimizing family.  Indeed convergence
in that compact family is strict perimeter convergence, hence its
perimeter measures converge weakly.  If the limit frontier meets the
open set $V$, its perimeter measure is positive on a nonnegative
compactly supported test in $V$, and the same holds for every sufficiently
late member of the sequence.  This proves openness.  The uniform density
bounds also give local Hausdorff convergence of the frontier supports:
a putative boundary point staying a positive distance from the limit
frontier would force positive volumes of both phases in a ball where
the limit is constant, contradicting local $L^1$ convergence.

The graph map
\[
 q_V:\mathcal H_V\longrightarrow C^\infty(B_1),\qquad q_V(F)=f_F,
\]
is continuous for the topology of smooth convergence on compact
subsets.  To see this, apply regular multiplicity-one convergence at
the limiting graph on each compact subball; the uniform graph bounds
and uniqueness identify every subsequential limit.  Define the finite
Borel measure
\[
 \lambda_V=(q_V)_*(\pi|_{\mathcal H_V})
\]
on the entire separable metrizable function space $C^\infty(B_1)$.
This definition identifies coincident local graphs without requiring
the image of $q_V$ to be a Borel subset of that space.  We write
$\mathcal G$ for this leaf-parameter space and $\lambda$ for
$\lambda_V$ when the box is fixed.

Every source contribution on $W$ is retained: a frontier meeting $W$
also meets $V$, and is therefore in $\mathcal H_V$.  No arbitrary
neighborhood of the reference minimizer is used to truncate this
measure.  Distinct realized graph parameters are disjoint, so for each
$y\in B_1$ the map
\begin{equation}\label{var:injective-leaf-evaluation}
 \ell\longmapsto f_\ell(y)
\end{equation}
is injective on the realized leaves.  All integrals with respect to
$\lambda$ mean the corresponding pushforward integrals over
$\mathcal H_V$; the injectivity assertion concerns the realized
graphs in those integrals.

\begin{lemma}[Atomic leaves]\label{var:atomic-leaves}
If $\lambda(\{\ell_0\})>0$, the corresponding regular sheet is
totally geodesic on $\operatorname{graph}(f_{\ell_0})\cap W$.
\end{lemma}

\begin{proof}
Fix any point of $\operatorname{graph}(f_{\ell_0})\cap W$ and take
signed normal-distance coordinates $(y,s)$ about that sheet on a
smaller patch compactly contained in $W$, writing the sheet as
$\Sigma=\{s=0\}$.  Put $m_0=\lambda(\{\ell_0\})>0$.
The $ss$ equation in \eqref{var:singular-system} is a one-dimensional
second-derivative equation after slicing by almost every normal line.
The bounded nonmeasure terms include $B$ and the coordinate connection
terms.  The atom at $s=0$ has coefficient $cm_0/2$: on the reference graph
the unit conormal is $ds$ and the sheet and volume Jacobians agree at
$s=0$.  No other graph produces an atom there, by
\eqref{var:injective-leaf-evaluation}.  Consequently $\partial_su$ has
bounded one-sided traces $v^\pm(y)$ for almost every $y$, and
\begin{equation}\label{var:atomic-normal-jump}
 v^+(y)-v^-(y)=\frac c2m_0.
\end{equation}
The existence and boundedness of these traces can also be read directly
from local semiconvexity of $u$ along the normal lines.

We compare the tangential equations from the two sides.  Let $A_s$ be the
second fundamental form of the parallel sheet $\{s=\mathrm{constant}\}$,
with normal $\partial_s$.  With $A_s(V,W)=g(\nabla_V\partial_s,W)$,
\begin{equation}\label{var:tangential-hessian-splitting}
 (\nabla^2u)|_{T\Sigma_s}
   =\nabla^2_{\Sigma_s}(u|_{\Sigma_s})+A_s\partial_su.
\end{equation}
Average this identity on each of $0<s<\varepsilon$ and
$-\varepsilon<s<0$, pairing with any smooth compactly supported tangential
tensor in the $y$ variables.  The intrinsic Hessian terms have the same
limit: integrate twice in $y$, then use uniform convergence of
$u(y,s)$ to $u(y,0)$ and smooth convergence of the parallel metrics.
The $B$ terms also have the same limit, since $u,X,\nabla X$ are
continuous.

The averaged tangential components of the sheet measure tend to zero.
Indeed the conormal of a nearby graph has tangential component bounded by
$C|s|$ relative to $\Sigma$ at a point of height $s$.  This is the
graph-difference estimate recalled above, with the smooth change from
graph coordinates to normal coordinates included.  Thus in a slab of
width $\varepsilon$ the tangential-tangential source is bounded in total
variation by $C\varepsilon^2$ times the total sheet measure in a fixed
box.  Dividing by $\varepsilon$ still gives a quantity tending to zero.
The atom on $\Sigma$ itself has no tangential component.

Subtracting the two limits in
\eqref{var:tangential-hessian-splitting} gives
$A_0(v^+-v^-)=0$ as a tensor distribution on $\Sigma$.
Equation~\eqref{var:atomic-normal-jump} and $c m_0>0$ imply $A_0=0$.
The sheet is smooth, so the conclusion holds at every point of the patch.
\end{proof}

\begin{lemma}[A common Jacobi field]\label{var:common-jacobi}
Let $D$ belong to the full-measure family above, and let $\Sigma$ be a
connected component of its regular interior frontier.  Suppose one of its
graph patches is a nonatom in the support of the corresponding leaf
measure.  Then there is a smooth function $\varphi>0$ on $\Sigma$ such
that
\begin{equation}\label{var:three-jacobi-equations}
 L_0\varphi=L_+\varphi=L_-\varphi=0.
\end{equation}
Here $L_0$, $L_+$, and $L_-$ are the outward normal linearizations of
$H$, $H+\operatorname{tr}_{\mathrm{tan}}K$, and
$H-\operatorname{tr}_{\mathrm{tan}}K$, respectively, at $\Sigma$.
All have principal part $-\Delta_\Sigma$, and
\begin{equation}\label{var:minimal-jacobi}
 L_0=-\Delta_\Sigma-
       \bigl(|A_\Sigma|^2+\operatorname{Ric}_g(\nu,\nu)\bigr).
\end{equation}
\end{lemma}

\begin{proof}
Choose a base point $p$ in the specified patch.  Since its graph
parameter is a nonatom in the support of the leaf measure, choose
distinct realized graphs converging smoothly on compact subpatches
to the reference graph, with representatives $D_j$ in the full-measure
family.  Such representatives exist because every graph neighborhood
has positive pushforward measure, whereas the reference fiber has
measure zero.  Compactness of $\mathcal F_*$ permits passing to a
subsequence with $D_j\to D'$ for some minimizing filled set $D'$.
We do not assert that $D'=D$ away from the selected component.

The limiting frontier of $D'$ contains the reference graph near $p$.
Indeed points on the converging graphs are frontier points of $D_j$,
and uniform density and local $L^1$ convergence put their limits on
the frontier of $D'$.  Set
\[
 C=\Sigma\cap\operatorname{frontier}(D').
\]
This set is nonempty and relatively closed in $\Sigma$.  It is also
relatively open: at each of its points the reference frontier is free
and regular, and the common-sheet property of
Proposition~\ref{var:area-envelope} makes the frontier of $D'$ agree
with it in a neighborhood.  Connectedness gives $C=\Sigma$.
In particular every point of $\Sigma$ is a regular point of the
limiting frontier $D'$.

Choose a nested exhaustion of $\Sigma$ by connected precompact smooth
domains $\Omega_m$ containing $p$, with
$\overline{\Omega_m}\subset\Omega_{m+1}$; if $\Sigma$ is compact, use
$\Sigma$ itself.  Use the fixed normal of
$\Sigma$ and its normal exponential map on a tubular neighborhood
of each compact stage.  Regular multiplicity-one convergence to
$D'$ gives, for sufficiently large $j$ at that stage, a unique normal
height $h_j$ over $\Omega_m$.  A finite covering of the slightly
larger compact stage by regular graph patches supplies this assertion
uniformly on $\overline{\Omega_m}$.  Uniqueness on each normal fiber
makes the heights agree on overlaps and on successive exhaustion
stages.  Thus there is no choice of branch or monodromy in this
construction.

The heights have a strict constant sign on each connected stage.
If $h_j$ vanished at a point, the common-sheet property would make
its zero set relatively open as well as closed there; connectedness
would force coincidence at $p$, contrary to the distinct initial
graphs.  Passing to a subsequence fixes the sign at $p$, and hence
on every exhaustion stage.  Replace $h_j$ by $-h_j$ in the linear
difference equations if that sign is negative, and normalize by
$h_j(p)>0$.  The equations $H=0$ and
$\operatorname{tr}_{\mathrm{tan}}K=0$ are unaffected by choosing
the fixed graph normal, even if a representative filled set induces
the opposite orientation on its local graph.
Subtract the three graph equations
\eqref{var:zero-sheet-expansions}.  The integral of the linearization
along the segment between the two graphs gives linear elliptic equations
for $h_j$, whose coefficients converge smoothly on compact subsets to
those of $L_0,L_+,L_-$.  Harnack's inequality and interior elliptic
estimates applied to the first equation bound $h_j/h_j(p)$ above and
below, with all derivatives, on each exhaustion stage.  A diagonal
subsequence converges smoothly to $\varphi$, with $\varphi(p)=1$.
Harnack's inequality gives positivity.  Passing to the limit in all
three equations proves \eqref{var:three-jacobi-equations}.  The same
argument applies when the approach is from the other side, since all
three linearized equations are homogeneous.
\end{proof}

The support qualification in this lemma loses no multiplier mass.  On
each graph box the complement of the support of $\lambda$ has zero
$\lambda$-measure.  A countable collection of such boxes suffices for
the regular incidence set, as explained below.

\subsection{Focusing with a measure-valued Hessian}

\begin{lemma}[Weak focusing]\label{var:weak-focusing}
On a free regular sheet satisfying
\eqref{var:zero-sheet-expansions}, define
\begin{equation}\label{var:Q-definitions}
 Q_+=u-g(X,\nu),\qquad Q_-=u+g(X,\nu).
\end{equation}
These functions are nonnegative and locally Lipschitz.  On the open part
of the sheet where $u>0$, they satisfy the distributional inequalities
\begin{equation}\label{var:weak-focusing-inequalities}
 L_\pm Q_\pm\leq -u|A_\Sigma\pm K_{\mathrm{tan}}|^2.
\end{equation}
\end{lemma}

\begin{proof}
We first identify the sign of the singular term in the stationary
curvature calculation.  For smooth positive $u$ and smooth $X$, set
\begin{equation}\label{var:local-stationary-metric}
 \mathbf g=-u^2 dz^2+
       g_{ij}(dx^i+X^i dz)(dx^j+X^j dz),\qquad
 \mathbf n=u^{-1}(\partial_z-X).
\end{equation}
The second fundamental form of a stationary slice is
$b=-\operatorname{sym}\nabla X/u$, with the convention for $K$ used in
the theorem.  Define
$\mathsf B(P,Q)=P:Q-(\operatorname{tr}P)(\operatorname{tr}Q)$.
Integrating the momentum term in the gravitational constraint pairing
at fixed lapse and shift gives the algebraic expression
\begin{align}\label{var:stationary-action-identity}
 u\bigl(R_g-\mathsf B(K,K)+2\mathsf B(K,b)\bigr)
 &=u\bigl(R_g+\mathsf B(b,b)
                    -\mathsf B(K-b,K-b)\bigr).
\end{align}
At $b=K$ its first metric variation agrees with that of the stationary
scalar-curvature action, up to divergences.  The spatial metric
coefficient of that variation is $-u\mathbf G_{ij}$, where $\mathbf G$
is the Einstein tensor of \eqref{var:local-stationary-metric}.
This identity can equally be obtained by substituting
$b=-\operatorname{sym}\nabla X/u$ and integrating the derivatives of a
compactly supported metric variation; it uses no field equation.

Let $\mathsf S$ be the Maxwell stress of a two-form whose slice electric
field is $c_nE$ and whose slice magnetic part is zero.  Only these slice
values are needed here.  Holding the electric flux form fixed under
$g'=h$ gives
\begin{equation}\label{var:fixed-flux-stress}
 \delta\bigl(2u c_n^2|E|^2\,dV_g\bigr)
 =u c_n^2(2E_iE_j-|E|^2g_{ij})h^{ij}\,dV_g
 =-u\mathsf S_{ij}h^{ij}\,dV_g.
\end{equation}
Adding the variation of the volume density in the full matter pairing
does not alter the adjoint coefficient, by
$u\mu_m+J_m(X)=0$.  Transport of a unit-ball ray contributes
$-(J_m)_{(i}X_{j)}$, where parentheses include the factor $1/2$.
It follows from the actual metric-adjoint expression in
\eqref{var:singular-system} that a Hessian source $M$ contributes
\begin{equation}\label{var:stress-trace-reversal}
 u(\mathbf G-\mathsf S)_{ij}
 =-(J_m)_{(i}X_{j)}+
       \bigl((\operatorname{tr}_g M)g-M\bigr)_{ij}.
\end{equation}
For a positive semidefinite $M$, the last tensor is positive
semidefinite: for a unit vector $v$, complete $v$ to an orthonormal
basis and sum $M$ on the remaining basis vectors.  In particular the
sign is preserved by trace reversal in precisely this direction.

The lower Hessian term can be read off explicitly.  Write
$\tau_b=\operatorname{tr}_g b$ and $(b^2)_{ij}=b_i{}^l b_{lj}$.
The spatial stationary Einstein identity is
\begin{align}
 u\mathbf G_{ij}
   &=(\Delta_g u)g_{ij}-(\nabla_g^2u)_{ij}
                      +(\mathcal C_b)_{ij},
 \label{var:explicit-spatial-einstein}\\
 \mathcal C_b
   &=u\left[\operatorname{Ric}_g+\tau_b b-2b^2
       -\tfrac12(R_g+\tau_b^2+|b|^2)g\right]
       -\mathcal L_Xb+X(\tau_b)g.
 \notag
\end{align}
At $b=K$, put
$T=u\mathsf S-\operatorname{sym}(J_m\otimes X^\flat)
       -\mathcal C_K$.
Equation~\eqref{var:stress-trace-reversal} then gives
\begin{equation}
 B(u,X,\nabla X)=\frac{\operatorname{tr}_gT}{n-1}g-T.
 \label{var:explicit-lower-hessian}
\end{equation}
Thus $B$ has no $du$ term.  All its coefficients are smooth
fixed-data coefficients multiplying $u,X,\nabla X$.
For the mollified pair, the same identity with $b=b_\epsilon$
shows that $C^1$ convergence of $b_\epsilon$ controls every
non-Hessian spatial term.

We justify the use of this calculation at the regularity of
\eqref{var:singular-system}.  Fix relatively compact coordinate patches
$U'\Subset U\Subset\{u>0\}$, and choose $a>0$ with $u\geq a$ on $U$.
All smoothing below is performed in $U$, then restricted to $U'$.
Let $\rho_\epsilon$ be a nonnegative standard mollifier, and convolve
the coordinate components of $u$ and $\beta=X^\flat$:
$u_\epsilon=u*\rho_\epsilon$,
$\beta_\epsilon=\beta*\rho_\epsilon$,
$X_\epsilon=g^{-1}\beta_\epsilon$.
The spatial metric is not changed.  In components the first equation
of \eqref{var:singular-system} reads
$\partial_{(i}\beta_{j)}=\Gamma_{ij}^{l}\beta_l-uK_{ij}$.
Hence the residual
\begin{equation}\label{var:KID-mollifier-residual}
 R_\epsilon:=\operatorname{sym}\nabla\beta_\epsilon+u_\epsilon K
 = (\Gamma\beta)*\rho_\epsilon-\Gamma\beta_\epsilon
       -(uK)*\rho_\epsilon+u_\epsilon K
\end{equation}
tends to zero in $C^1(U')$.  More explicitly, for a smooth coefficient
$A$ and a Lipschitz function $f$,
\begin{equation}\label{var:C1-commutator}
 \|(Af)*\rho_\epsilon-A(f*\rho_\epsilon)\|_{C^1(U')}
 \leq C\epsilon\bigl(
       \|A\|_{C^2(U)}\|f\|_{L^\infty(U)}
       +\|A\|_{C^1(U)}\|Df\|_{L^\infty(U)}\bigr).
\end{equation}
For its derivative, differentiate the product to obtain the sum of
the commutator of $DA$ with $f$ and the commutator of $A$ with $Df$.
Both are bounded by the first moment of $\rho_\epsilon$ times the
displayed norms.  No second derivative of $f$ is used.
Since $u_\epsilon\geq a/2$ and $Du_\epsilon$ is uniformly bounded,
\begin{equation}\label{var:b-mollifier-convergence}
 b_\epsilon:=-\operatorname{sym}\nabla X_\epsilon/u_\epsilon
       =K-R_\epsilon/u_\epsilon\longrightarrow K
       \quad\hbox{in }C^1(U').
\end{equation}

In the Hessian equation, first express the measure-valued components
relative to coordinate Lebesgue measure, dividing by the smooth
positive volume density.  Positivity is unchanged.  Convolve this
already un-trace-reversed equation.  It gives
\begin{equation}\label{var:Hessian-mollifier}
 \nabla_g^2u_\epsilon
 =B_\epsilon+\mathcal M_\epsilon+
        (\Gamma\,du)*\rho_\epsilon-\Gamma\,du_\epsilon,
 \qquad B_\epsilon\longrightarrow B\quad\hbox{uniformly on }U'.
\end{equation}
The connection commutator is $O(\epsilon)$ in $C^0$ because $du$ is
bounded.  For every $x$, $\mathcal M_\epsilon(x)$ is a positive
semidefinite matrix.  We trace this matrix using $g(x)$ only after
convolution.  Consequently
\begin{equation}\label{var:exact-PSD-reversal}
 P_\epsilon(x):=
       (\operatorname{tr}_{g(x)}\mathcal M_\epsilon(x))g(x)
                        -\mathcal M_\epsilon(x)\geq0
\end{equation}
exactly, even if its norm diverges as $\epsilon\downarrow0$.
There is no product commutator involving a variable metric times an
unbounded singular Hessian.

Apply the smooth stationary calculation to $u_\epsilon,X_\epsilon$
and $b_\epsilon$.  Formula~\eqref{var:b-mollifier-convergence} controls
the induced constraint densities and all the non-Hessian spatial
adjoint terms uniformly.  Indeed \eqref{var:explicit-spatial-einstein} uses only
$b_\epsilon,\nabla b_\epsilon,u_\epsilon,X_\epsilon,
\nabla X_\epsilon$ outside the Hessian term.
Equation~\eqref{var:explicit-lower-hessian} shows that
uniform convergence of $du_\epsilon$ is unnecessary.
Writing $\mathbf T_\epsilon=\mathbf G_\epsilon-\mathsf S_\epsilon$,
we obtain uniformly on $U'$
\begin{align}\label{var:smoothed-stress-blocks}
 \mathbf T_\epsilon(\mathbf n_\epsilon,\mathbf n_\epsilon)
     &=\mu_m+o(1),&
 \mathbf T_\epsilon(\mathbf n_\epsilon,\cdot)|_{TU'}
     &=J_m+o(1),\\
 u_\epsilon(\mathbf T_\epsilon)_{ij}
     &=-(J_m)_{(i}X_{j)}+(P_\epsilon)_{ij}+o(1).&&
 \notag
\end{align}
To see the null sign of the bounded terms, when $\mu_m>0$,
\eqref{var:moment-complementarity} forces
$|X|=u$ and $J_m=-\mu_m X^\flat/u$.  Their limiting block tensor is
$\mu_m\ell\otimes\ell$, with
$\ell(\mathbf n)=1$ and $\ell(V)=-g(X,V)/u$.
When $\mu_m=0$, all its blocks vanish.  This description, or the
direct formula $u\mu_m(1\mp g(X,\nu)/u)^2$, proves nonnegativity on
$\mathbf n\pm\nu$.  Maxwell stress is also nonnegative on null
vectors.  The exact positivity of $P_\epsilon$ therefore gives
\begin{equation}\label{var:mollified-null-Ricci}
 u_\epsilon\operatorname{Ric}_{\mathbf g_\epsilon}
       (\mathbf n_\epsilon\pm\nu,
        \mathbf n_\epsilon\pm\nu)\geq-o(1).
\end{equation}
This is a one-sided estimate; it requires no uniform bound on the full
spacetime curvature.

For completeness, apply the null shape-operator equation to the
outgoing null geodesics orthogonal to the fixed sheet in the smooth
metric $\mathbf g_\epsilon$.  Use nearby stationary slices as the
parameter surfaces.  At the initial slice the null velocity with
$dz/d\tau=1$ is $u_\epsilon(\mathbf n_\epsilon+\nu)$.
Subtracting the stationary translation $\partial_z$ leaves the
spatial velocity $u_\epsilon\nu-X_\epsilon$, whose normal component
is $u_\epsilon-g(X_\epsilon,\nu)$.  The Riccati equation gives the
expansion derivative
\begin{equation}\label{var:smoothed-focusing-formula}
 -u_\epsilon\bigl(
     |A_\Sigma+(b_\epsilon)_{\mathrm{tan}}|^2+
     \operatorname{Ric}_{\mathbf g_\epsilon}
        (\mathbf n_\epsilon+\nu,\mathbf n_\epsilon+\nu)\bigr),
\end{equation}
up to a bounded renormalization factor times the initial expansion.
Tangential transport contributes a bounded tangential velocity times
the tangential derivative of that expansion.  Both error terms tend
to zero uniformly: $H+\operatorname{tr}_\Sigma K=0$ and
$b_\epsilon\to K$ in $C^1$.
The spatial expansion linearization has coefficients converging
uniformly to those of $L_+$.  Its principal part is always the fixed
$-\Delta_\Sigma$; convergence of its action on
$u_\epsilon-g(X_\epsilon,\nu)$ therefore holds in distributions,
using uniform convergence of these functions and uniform bounds on
their first derivatives.  Combining
\eqref{var:mollified-null-Ricci}--\eqref{var:smoothed-focusing-formula}
proves the plus inequality in
\eqref{var:weak-focusing-inequalities}.  Replace the future normal by
$-\mathbf n_\epsilon$, equivalently reverse stationary time and
$X_\epsilon,b_\epsilon$, to obtain the minus inequality.  Since
$U'\Subset\{u>0\}$ was arbitrary, this proves the stated local result.
\end{proof}

\subsection{Positive Jacobi fields at singular ends}

We give the analytic details needed at singular points.  Their purpose
is stronger than removability of a bounded subsolution: a bounded
reciprocal Jacobi field must tend to zero at every singular end.

\begin{lemma}[Capacity and a local mean-value estimate]
\label{var:capacity-submean}
Let $T$ be a multiplicity-one locally perimeter-minimizing hypersurface
in a compact smooth ambient patch, of dimension $k\geq3$.  Its singular
set has zero surface $2$-capacity.  A bounded nonnegative continuous function $w$
on the regular part, satisfying $\Delta_T w\geq-C_0$ there, extends
weakly across the singular set and satisfies, for sufficiently small
ambient balls centered on $\operatorname{supp}T$,
\begin{equation}\label{var:surface-submean}
 \sup_{T\cap B_{r/2}}w
 \leq C\left[
       \left(r^{-k}\int_{T\cap B_r}w^2\,dA\right)^{1/2}
          +C_0r^2\right].
\end{equation}
The supremum is over regular points.  Constants are uniform in a
fixed compact patch, and under its sufficiently small rescalings.
The assertion also holds for a function defined on one regular
component and extended by zero on the other components.
\end{lemma}

\begin{proof}
The singular dimension bound is $k-7$, and local perimeter bounds give
$\Area(T\cap B_\rho)\leq C\rho^k$; see
\cite{SimonGMT2014}.  If the singular set is nonempty, cover a compact
part of it by balls $B_{r_i}$ with $r_i$ arbitrarily small and
$\sum_i r_i^{k-2}$ arbitrarily small.  Choose cutoffs equal to one on
$B_{r_i}$, supported on $B_{2r_i}$, with gradient bounded by $C/r_i$.
The maximum of these cutoffs has Dirichlet energy at most
$C\sum_i r_i^{k-2}$ and support of area tending to zero.  Truncation
and smooth approximation produce cutoffs $\chi_j$, vanishing near
the singular set and tending to one almost everywhere, with
\begin{equation}\label{var:capacity-cutoffs}
 \int|\nabla_T\chi_j|^2\,dA\longrightarrow0.
\end{equation}
The same construction applies on compact cone annuli and their
spherical links, where the singular link dimension is at most $k-8$.

On the regular part, test the subsolution inequality by
$\eta^2\chi_j^2w$, using smooth approximations or positive
truncations if needed.  Integration by parts and Young's inequality
bound the Dirichlet integral of $w$ on smaller patches by
\begin{equation}\label{var:subsolution-Caccioppoli}
 \int\eta^2\chi_j^2|\nabla_Tw|^2
 \leq C\|w\|_\infty^2
       \int\bigl(|\nabla_T\eta|^2+
                     \eta^2|\nabla_T\chi_j|^2\bigr)
       +C C_0\|w\|_\infty\int\eta^2.
\end{equation}
Thus $w$ has locally square-integrable weak gradient across the
singular set.  In a bounded compact test, the extra terms containing
$\nabla\chi_j$ tend to zero by Cauchy--Schwarz and
\eqref{var:capacity-cutoffs}.  This proves the weak extension.
The same proof on any union of regular components is valid, since
its boundary within the support lies in the singular set.

The Michael--Simon surface Sobolev inequality
\cite{MichaelSimon1973}, applied in coordinate or isometrically
embedded compact patches, gives for compactly supported $v$
\begin{equation}\label{var:surface-Sobolev}
 \left(\int|v|^{2k/(k-2)}\,dA\right)^{(k-2)/k}
 \leq C\int\bigl(|\nabla_Tv|^2+v^2\bigr)\,dA.
\end{equation}
The ambient embedding contributes a uniformly bounded mean-curvature
term.  The inequality extends to the above weak functions by the
same capacity cutoffs.  After scaling a ball to radius one, apply
the energy estimate to successive positive powers of
$w+C_0r^2$ and nested radial cutoffs.  The Sobolev exponent multiplies
each power by $k/(k-2)$; the resulting geometric iteration gives
the $L^2$-to-$L^\infty$ estimate
\eqref{var:surface-submean}.  All constants depend only on the
ambient patch, dimension, and perimeter bound.  Their scaled versions
remain bounded as $r\downarrow0$.
\end{proof}

\begin{lemma}[A bounded cone subsolution vanishes]
\label{var:cone-subsolution}
Let $C\subset\mathbb R^{k+1}$ be a nonflat, multiplicity-one,
perimeter-minimizing boundary cone.  A bounded nonnegative continuous function
$z$ on its entire regular part satisfying
\begin{equation}\label{var:cone-subsolution-equation}
 \Delta_C z\geq |A_C|^2z
\end{equation}
in distributions is identically zero.  The function may be zero on
some regular components.
\end{lemma}

\begin{proof}
Only $k\geq7$ can occur.  Let $L=C\cap\mathbb S^k$ be the link,
let $s=\log r$, and put $V=|A_L|^2$.  The inequality becomes
\begin{equation}\label{var:log-radius-inequality}
 z_{ss}+(k-2)z_s+\Delta_Lz\geq Vz.
\end{equation}
We first justify testing across the singular link.  On each compact
axial slab, the Caccioppoli test $\eta^2\chi_j^2z$ gives local
$W^{1,2}$ bounds and a bound for
$\int V\eta^2\chi_j^2 z^2$; the axial first-order term is bounded
by the same energy inequality.  Passing through the capacity
cutoffs of Lemma~\ref{var:capacity-submean} proves the weak inequality
across the singular set.  Testing with a nonnegative function equal
to one on a smaller slab also bounds $\int Vz$ there.  Positivity
allows Fatou's lemma for this term.  The terms involving
$\nabla\chi_j$ tend to zero by the energy bound and
\eqref{var:capacity-cutoffs}.  These observations justify integration
of \eqref{var:log-radius-inequality} over the whole link against any
nonnegative compactly supported axial test.

In particular the bounded nonnegative function
$q(s)=\int_Lz(s,\theta)\,dA_L(\theta)$ satisfies
\begin{equation}\label{var:log-radius-monotonicity}
 q''+(k-2)q'\geq0
\end{equation}
on the whole real line.  Convolve in $s$.  Each smooth bounded
convolution satisfies the same inequality, so
$e^{(k-2)s}q'(s)$ is nondecreasing for that convolution.  If
$q'(s_0)<0$, then for $s<s_0$
$q'(s)\leq q'(s_0)e^{(k-2)(s_0-s)}$.
Integrating backward contradicts boundedness.  Thus every
convolution is nondecreasing, and the original $q$ has a
nondecreasing representative.  If $z$ is nonzero, its forward limit
$q_\infty$ is strictly positive.

Choose $w\in C_c^\infty((1,2))$, $w\geq0$, with integral one, and
put $w_T(s)=T^{-1}w(s/T)$.  The functions
\begin{equation}\label{var:forward-average}
 \overline z_T(\theta)=\int_{\mathbb R}w_T(s)z(s,\theta)\,ds
\end{equation}
are bounded uniformly, and their integrals tend to $q_\infty$.
Take a weak-star convergent subsequence in $L^\infty(L)$ with limit
$\overline z\geq0$.  For a compactly supported nonnegative regular
link test $\psi$, integration by parts in $s$ bounds the axial terms by
$\|z\|_\infty\|\psi\|_{L^1(L)}(\|w_T''\|_1+(k-2)\|w_T'\|_1)$, which tends to zero.
Therefore
\begin{equation}\label{var:averaged-link-inequality}
 \Delta_L\overline z\geq V\overline z,
 \qquad \int_L\overline z\,dA_L=q_\infty>0.
\end{equation}
Bounded subsolution regularity and the capacity argument give the
needed local weak derivatives.  Test this inequality by
$\chi_j^2\overline z$ and use Young's inequality to obtain
\begin{equation}\label{var:link-energy-vanishing}
 \frac12\int_L\chi_j^2|\nabla_L\overline z|^2
       +\int_LV\chi_j^2\overline z^2
 \leq2\|\overline z\|_\infty^2
                  \int_L|\nabla_L\chi_j|^2\longrightarrow0.
\end{equation}
Thus $\overline z$ is constant on every connected regular component
of $L$, and a positive constant can occur only on a totally geodesic
component.

We show that a nonflat minimizing boundary cone has no such
component.  A connected totally geodesic regular component lies in
an equator $Q\cong\mathbb S^{k-1}$.  It is relatively open there;
its relative boundary lies in the singular set, since at a regular
boundary point the same smooth sheet continues.  A closed set of
Hausdorff dimension less than $(k-1)-1$ does not disconnect $Q$:
in coordinate balls one connects points by polygonal paths whose
vertices are chosen outside the measure-zero exceptional sets of
segments meeting that set, and then concatenates such paths.
Here the singular set has dimension at most $k-8$.
It follows that this component contains $Q$ minus the singular set,
and its closure contains the entire equator.

Any other regular component lies strictly in one open hemisphere
determined by $Q$.  It cannot meet $Q$ at a regular point without
being part of the same regular sheet.  On that component the signed
coordinate height $h$ is strictly of one sign and satisfies
$\Delta_Lh=-(k-1)h$.  It and its gradient are bounded.  Integrate
against capacity cutoffs, extended by zero on other components.
The boundary terms tend to zero by Cauchy--Schwarz; hence
$\int h=0$, a contradiction.  There are no other regular
components.  Regular points are dense in the support, and
multiplicity one now says that $C$ is the plane through $Q$.
The componentwise interpretation is also consistent with the
decomposition of a multiplicity-one oriented minimizing boundary
into minimizing regular components; see
\cite[Chapter~7, Section~5, Corollary~5.11 and Remark~5.12]{SimonGMT2014}.
This contradicts nonflatness.  Therefore
\eqref{var:averaged-link-inequality} is impossible and $z=0$.
\end{proof}

\begin{lemma}[Reciprocal Jacobi decay]\label{var:reciprocal-jacobi}
Let $\Sigma$ be a connected regular component of a compact interior
minimizing frontier, and let $\varphi>0$ solve
\eqref{var:minimal-jacobi}.  Then
\begin{equation}\label{var:Jacobi-singular-behavior}
 \inf_\Sigma\varphi>0,
 \qquad
 \varphi(x)\longrightarrow+\infty
 \quad\hbox{as }x\in\Sigma\hbox{ approaches the singular set}.
\end{equation}
The second assertion is vacuous for a smooth compact component.
\end{lemma}

\begin{proof}
The Ricci curvature is bounded on a compact neighborhood of the
frontier; choose $C$ with $\operatorname{Ric}_g(\nu,\nu)\geq-C$.
For $\delta>0$ put $w_\delta=(1-\varphi/\delta)_+$ on $\Sigma$,
and extend it by zero on other regular components.  Convex
truncation of the Jacobi equation gives
$\Delta w_\delta\geq-C$; on $\{\varphi<\delta\}$ this follows
from $-\Delta\varphi/\delta=(|A|^2+\operatorname{Ric}(\nu,\nu))
\varphi/\delta\geq-C$, and the truncation contributes a nonnegative
measure.  Moreover $0\leq w_\delta\leq1$ and $w_\delta\to0$
almost everywhere as $\delta\downarrow0$.

Choose first a fixed radius $r>0$ so small that the $Cr^2$ term
in \eqref{var:surface-submean}, including its constant, is less
than $1/4$.  Cover the compact frontier by finitely many balls
of radius $r/2$.  Dominated convergence in their radius-$r$
enlargements then permits choosing $\delta>0$ so that the integral
term is less than $1/4$ in every ball.  Thus $w_\delta\leq1/2$
on $\Sigma$, and $\varphi\geq\delta/2$.  This proves the first
assertion with the required order of choices, $r$ before $\delta$.

Set $z=1/\varphi$ and extend it by zero to the other regular
components.  It is bounded and satisfies
\begin{equation}\label{var:reciprocal-inequality}
 \Delta_\Sigma z
   =\bigl(|A|^2+\operatorname{Ric}_g(\nu,\nu)\bigr)z
                  +2\varphi^{-3}|\nabla_\Sigma\varphi|^2
   \geq(|A|^2-C)z.
\end{equation}
Fix a singular point $p$ and any sequence of scales $r_j\downarrow0$.
Pass to a subsequence along which the rescaled frontiers converge
to a multiplicity-one minimizing tangent cone $C_p$.  The cone
is nonflat: a multiplicity-one planar tangent cone would imply
regularity of the frontier at $p$.  These are the standard
compactness and excess-regularity conclusions for minimizing
boundaries \cite{SimonGMT2014}.

On compact subsets of the regular cone the convergence is smooth.
The rescaled $z$ have subsequences converging locally uniformly to
a bounded nonnegative function $z_0$.  Here is a useful precise
way to obtain this convergence.  On a connected regular patch,
either $z$ tends to zero at one base point, in which case Harnack
for the positive Jacobi function $\varphi$ forces $z\to0$ on
smaller patches, or $z$ stays bounded below there along a
subsequence.  In the latter case the same Harnack inequality and
interior estimates bound $\varphi$ and its derivatives, giving
smooth convergence of $z$.  The uniform upper bound for $z$
was already established.  Branches on which the zero extension
is used cause no difficulty; choose the branch in each regular
patch and diagonalize.  Thus one obtains $z_0$ on the whole
regular cone, allowing zero components.
Scaling \eqref{var:reciprocal-inequality} and passing to regular
patches gives $\Delta_{C_p}z_0\geq|A_{C_p}|^2z_0$.
Lemma~\ref{var:cone-subsolution} implies $z_0=0$.

This regular convergence implies vanishing of the normalized
$L^2$ integral of $z$ on each fixed bounded rescaled ball.
To check points near the singular cone, first discard a small
neighborhood of that set and of the cone vertex.  Their limiting
area can be made arbitrarily small.  Area convergence and the
uniform bound for $z$ control the discarded part, while smooth
regular convergence controls its complement.  The scaled
version of \eqref{var:surface-submean}, with source bounded by
$Cr_j^2\|z\|_\infty$, now gives uniform decay on smaller balls.

In particular, if $x_j\in\Sigma$ tends to $p$, choose scales
$r_j$ comparable to $d_g(x_j,p)$.  The preceding argument and
the mean-value estimate show $z(x_j)\to0$; any contrary
subsequence would have a tangent-cone subsequence of the type
just treated.  A sequence approaching the compact singular set
has a subsequence tending to a particular singular point.
Applying this fixed-point argument to that subsequence proves
the uniform sequential assertion in
\eqref{var:Jacobi-singular-behavior}.
\end{proof}

\subsection{Diffuse mass at a zero lapse and the final curvature bound}

\begin{lemma}[No diffuse leaf mass on the zero set]
\label{var:diffuse-zero-mass}
In a uniform graph box, let $\lambda^{\mathrm{na}}$ denote the
nonatomic part of its leaf measure.  Let
\[
 \mathcal Z_W=\{\ell:
       \operatorname{graph}(f_\ell)\cap W\ne\varnothing,
       \ u=0\text{ on }\operatorname{graph}(f_\ell)\cap W\}.
\]
Then
\begin{equation}\label{var:diffuse-zero-claim}
 \lambda^{\mathrm{na}}(\mathcal Z_W)=0.
\end{equation}
\end{lemma}

\begin{proof}
On $W=B_0\times I_0$, the coordinate $tt$ equation in
\eqref{var:singular-system} is the following equality of measures:
\begin{equation}\label{var:graph-disintegration}
 \partial_t^2u
 =b(y,t)\,dy\,dt+
   dy\int_{\mathcal G}w(\ell,y)
       \mathbf1_{I_0}(f_\ell(y))
       \delta_{f_\ell(y)}(dt)\,d\lambda(\ell).
\end{equation}
Here $y\in B_0$ and the left side and bounded-density term are
restricted to $W$ as measures.  The indicator records the restriction
of the graph measure to the working $t$-interval; it is not a
differentiation of a characteristic function.  The density $b$ is
bounded, and $0<w_0\leq w(\ell,y)\leq w_1$ for the contributing
graphs.  The weight is $c/2$ times the squared conormal on
$\partial_t$, times the graph area Jacobian, divided by the volume
Jacobian.  All contributions on $W$ are included because their
frontiers lie in $\mathcal H_V$.  Pairing with product tests and
using the graph area formula proves the display; measure
approximation then gives its linewise disintegration.

For almost every $y$, $t\mapsto u(y,t)$ is Lipschitz and its
derivative $v_y=\partial_tu(y,\cdot)$ is a one-dimensional BV
function.  This follows from the sliced equation, or from
semiconvexity and Fubini.  Write
$Z_y=\{t:u(y,t)=0\}$.  At any point of $Z_y$ where the two
one-sided precise values of $v_y$ agree, their common value is
zero: a nonzero common value would contradict the nonnegative
minimum of the primitive $u(y,\cdot)$ at that point.
The diffuse derivative of a one-dimensional BV function does
not charge a fixed precise-value level.  In the usual
$D^av+D^cv$ notation this locality statement is
\begin{equation}\label{var:BV-level-locality}
 |D^av+D^cv|\bigl(\{\widetilde v=0\}\bigr)=0;
\end{equation}
it follows from the BV chain rule, or directly from the
one-dimensional coarea formula; see
\cite[Proposition~8, equation~(16)]{AmbrosioDeLellisMaly2007}.
Consequently the diffuse part of $Dv_y$ gives no mass to $Z_y$.
The countably many jump points do not affect that assertion.

The bounded term $b$ is zero almost everywhere on $\{u=0\}$.
Indeed $X=0$ there by $u\geq|X|$, and the first derivatives of
$u$ and $X$ vanish almost everywhere on their respective zero
sets.  Thus the linear expression $B(u,X,\nabla X)$ and the
connection term involving $du$ both vanish almost everywhere
there.  Fubini gives $b(y,\cdot)=0$ almost everywhere on $Z_y$
for almost every line.

Finally, by \eqref{var:injective-leaf-evaluation} the pushforward
of $w(\ell,y)d\lambda^{\mathrm{na}}(\ell)$ under
$\ell\mapsto f_\ell(y)$ is atomless for every $y$.
All line atoms of the area source in
\eqref{var:graph-disintegration} therefore come from transverse
atoms of $\lambda$; there are no hidden atoms produced by the
evaluation map.  On almost every line, restrict the diffuse part
of \eqref{var:graph-disintegration} to $Z_y$ and use
\eqref{var:BV-level-locality} and the preceding vanishing of $b$.
It follows that
\begin{equation}\label{var:linewise-zero-mass}
 \int_{\mathcal G}w(\ell,y)
       \mathbf1_{I_0}(f_\ell(y))
       \mathbf1_{\{u(y,f_\ell(y))=0\}}
                  \,d\lambda^{\mathrm{na}}(\ell)=0
\end{equation}
for almost every $y$.  The integrand is defined to be zero whenever
$f_\ell(y)\notin I_0$, so $u$ is evaluated only inside $W$.
For each leaf set
\[
 m_W(\ell)=\int_{B_0}w(\ell,y)
                    \mathbf1_{I_0}(f_\ell(y))\,dy.
\]
This quantity is strictly positive whenever its graph meets $W$:
that intersection projects to a nonempty open subset of $B_0$, and
$w\geq w_0>0$.  A uniform positive lower bound for $m_W$ over all
leaves is neither asserted nor needed.  Integrate
\eqref{var:linewise-zero-mass} in $y$ and use Tonelli's theorem.
On $\mathcal Z_W$ the result is
\[
 \int_{\mathcal Z_W}m_W(\ell)\,d\lambda^{\mathrm{na}}(\ell)=0.
\]
Strict positivity of $m_W$ there gives
$\lambda^{\mathrm{na}}(\mathcal Z_W)=0$, proving the claim.
\end{proof}

\begin{proposition}[Elimination of the interior area multiplier]
\label{var:area-elimination}
Under the hypotheses above and the exclusion of a smooth interior
enclosing frontier with zero future expansion in
Definition~\ref{def:rigidity-class}, the measure $\pi$ is
concentrated on the given boundary:
\begin{equation}\label{var:pi-boundary-only}
 \pi=\delta_S,
 \qquad
 \int_{\mathcal F_*}a'_D(h)\,d\pi(D)=a'_S(h).
\end{equation}
Here $\delta_S$ denotes the point mass at the filled minimizer
whose frontier is $S$.
\end{proposition}

\begin{proof}
First consider a connected regular component $\Sigma$ with a
nonatomic support patch.  Choose the common field $\varphi$ of
Lemma~\ref{var:common-jacobi}.  The functions $Q_\pm$ of
\eqref{var:Q-definitions} are bounded on the compact frontier.
By Lemma~\ref{var:reciprocal-jacobi},
$Q_\pm/\varphi$ tends to zero at every singular end.  If one of
these ratios is positive somewhere, it therefore attains a
positive maximum at a regular interior point.  This is also
immediate when $\Sigma$ is smooth and compact.

Write $Q_\pm=\varphi v$.  For an operator
$L=-\Delta+b\cdot\nabla+d$ with $L\varphi=0$, direct product
differentiation gives
\begin{equation}\label{var:ground-state-transform}
 L(\varphi v)=\varphi\left[
   -\Delta v+(b-2\nabla\log\varphi)\cdot\nabla v\right].
\end{equation}
At any point where $Q_\pm>0$, one has $u>0$, so
Lemma~\ref{var:weak-focusing} applies.  The weak strong maximum
principle for the locally Lipschitz $v$ and the smooth coefficients
in \eqref{var:ground-state-transform} implies that $v$ is constant
at its positive maximum on the connected component of
$\{Q_\pm>0\}$ containing that point.  This component is also
closed in $\Sigma$: at a regular limit point,
$Q_\pm=v\varphi>0$ by continuity.  Since $\Sigma$ is connected,
it is all of $\Sigma$.  Substitution in
\eqref{var:weak-focusing-inequalities} then gives
\begin{equation}\label{var:sheet-curvature-bound}
 A_\Sigma\pm K_{\mathrm{tan}}=0,
 \qquad |A_\Sigma|\leq |K|.
\end{equation}
The weak maximum principle used here is the ordinary local
elliptic one, for example \cite{GilbargTrudinger2001}; no
maximum principle on an unresolved singular boundary is used.
If neither ratio is positive anywhere, $Q_+=Q_-=0$ on
$\Sigma$, hence $u=0$ there and $X=0$ there by domination.
Lemma~\ref{var:diffuse-zero-mass} excludes multiplier mass on
this alternative in every nonatomic patch.  In atomic patches
Lemma~\ref{var:atomic-leaves} already gives $A=0$.

We make the exceptional set independent of the regular incidence point.
The open inner boxes $W$ constructed above cover all free regular
points of all minimizing frontiers.  Their union is an open subset
of the second-countable ambient manifold, so choose a countable
subcover $W_i$, retaining the corresponding middle boxes $V_i$ and
graph maps $q_i$ on the open families $\mathcal H_{V_i}$.  In
particular the open sets $\mathcal H_{V_i}\times W_i$ cover the
regular incidence set in the family-times-manifold topology.

For each $i$ let $\lambda_i=(q_i)_*(\pi|_{\mathcal H_{V_i}})$,
let $\mathcal A_i$ be its countable set of atoms, and let
$\mathcal Z_i=\mathcal Z_{W_i}$ be the zero-leaf set from
Lemma~\ref{var:diffuse-zero-mass}.  These are measurable sets.
For the last assertion, the condition that a graph meets $W_i$
is open in the graph topology.  The condition that $u$ vanishes
on its intersection with $W_i$ can be checked at a countable
dense set of base points: failure gives, by continuity, an open
base set on which the graph stays inside $W_i$ and $u>0$.
Thus it is a Borel condition.

The set
\[
 \mathcal N_i=
    \bigl(C^\infty(B_{1,i})\setminus\operatorname{supp}\lambda_i\bigr)
         \cup(\mathcal Z_i\setminus\mathcal A_i)
\]
has $\lambda_i$-measure zero, since the second term has
$\lambda_i$-measure $\lambda_i^{\mathrm{na}}(\mathcal Z_i)=0$.
Consequently
\[
 \pi\bigl(\{F\in\mathcal H_{V_i}:q_i(F)\in\mathcal N_i\}\bigr)=0.
\]
Remove the countable union of these pullbacks, together with the
already excluded null set in the good family.  For a remaining
frontier and any of its regular points $p$, choose $i$ with
$p\in W_i$.  Its leaf is either an atom, giving $A=0$ locally,
or a nonatom in the support of $\lambda_i$.  In the latter case
the common Jacobi argument applies on its entire connected
regular component.  If that component had $u=0$ everywhere,
the leaf would belong to $\mathcal Z_i\setminus\mathcal A_i$,
which was removed.  Therefore the positive-ratio alternative
gives $|A|\leq |K|$.  This proves the curvature bound
simultaneously at every regular point of every remaining
frontier.
In particular,
\begin{equation}\label{var:uniform-frontier-curvature}
 |A|\leq\sup_{\mathcal K}|K|,
\end{equation}
where $\mathcal K$ is a common compact set containing the
minimizing frontiers.

Fix such an interior frontier and any proposed singular point.
Blow it up at scales tending to zero.  The tangent cone is a
multiplicity-one minimizing boundary cone, with smooth
convergence on its regular part.  The rescaled bound
\eqref{var:uniform-frontier-curvature} tends to zero, so every
regular component of the cone is totally geodesic.  The
equator and height argument in
Lemma~\ref{var:cone-subsolution} shows that the cone is a plane.
Multiplicity-one planar tangent regularity contradicts the
point being singular.  The entire compact frontier is therefore
smooth.

This frontier is a genuine enclosing cut.  In particular the
minimizing filled set has no removable bounded open component
of its complement: filling one would preserve the obstacle
condition and strictly reduce its full perimeter.  Its
components carry the outward normals of that enclosing set,
and \eqref{var:zero-sheet-expansions} gives zero future expansion
on every component.  The assumed exclusion of a smooth
interior enclosing marginal frontier rules it out.
Consequently $\pi$-almost every frontier is $S$.  The filled
set with frontier $S$ and the prescribed obstacle is unique
modulo null sets, and $\pi$ is a probability.  This proves
\eqref{var:pi-boundary-only}.
\end{proof}
\subsection{Smooth stationarity and the boundary laws}

We finish the original-data variation argument.  The lapse $u$ and
shift $X$ in the following proposition are the measure moments
constructed above; they are unrelated to the lapse notation used
locally in the filled scalar construction.

\begin{proposition}[Stationary data supplied by equality]
\label{var:stationary-data}
\label{var:stationary-development}
There are smooth fields $u,X$ on $\Omega$, including $S$, with
$u\ge |X|$, $u>0$ in the open exterior and
\begin{equation}
 \operatorname{sym}\nabla X=-uK,\qquad
 u\mu_m+J_m(X)=0.
 \label{var:smooth-kid-complementarity}
\end{equation}
For any sufficiently small decrease $q_0<q$ with
$q_0>(n-2)/2$, their end behavior is
\begin{equation}
 (u,X)=(b^0,b^i)+O_2(r^{-q_0}).
 \label{var:lapse-shift-end}
\end{equation}
On $\R_z\times\operatorname{int}\Omega$ set
\begin{equation}
 \mathbf g=-u^2\dd z^2+
             g_{ij}(\dd x^i+X^i\dd z)(\dd x^j+X^j\dd z),
 \qquad \xi=\partial_z,\qquad
 \mathbf F=c_n\dd z\wedge uE^\flat.
 \label{var:stationary-product}
\end{equation}
The stationary fields satisfy Maxwell's equations and
\begin{equation}
 \operatorname{Ein}_{\mathbf g}-\mathsf S
       =\mu_m u^{-2}\xi^\flat\otimes\xi^\flat,
 \qquad
 \mathsf S_{ab}=2\left(\mathbf F_{ac}\mathbf F_b{}^c
           -\tfrac14\mathbf F_{cd}\mathbf F^{cd}\mathbf g_{ab}\right).
 \label{var:stationary-einstein}
\end{equation}
The extra term is zero unless $\xi$ is null.  The original slice
has future normal $N=(\partial_z-X)/u$ and induces exactly
$(g,K,E)$, with $\mathbf F$ restricting to zero on it.
There is a global smooth function $p$, including up to $S$, with
\begin{equation}
 \dd p=(k-1)uE^\flat,
 \qquad p(\infty)=0,
 \qquad p=O_2(r^{2-n}),
 \qquad \dd(i_X\alpha)=0.
 \label{var:potential-maxwell}
\end{equation}

Let $\kappa=c/2>0$.  On $S$ the boundary laws are
\begin{equation}
 X=u\nu,\qquad
 \partial_\nu u+K(X,\nu)=\kappa,\qquad
 cH=2L_+^*u,
 \label{var:horizon-boundary-laws}
\end{equation}
where $L_+$ is the outward normal linearization of
$H+\operatorname{tr}_S K$, and its adjoint uses $\dd A_g$.
Either $u>0$ on all of $S$, or $u=0$ on all of $S$.
The nonnegative function $W=u^2-|X|^2$ tends to one at
infinity, vanishes on $S$, and is strictly positive in a
deleted collar of $S$.  If $u>0$ on $S$, then
\begin{equation}
 \dd W=2\kappa X^\flat\quad\hbox{on }S.
 \label{var:positive-lapse-boundary}
\end{equation}
If $u=0$ on $S$, writing $s$ for exterior distance from $S$,
there are smooth fields $a_*,Y_*$ such that
\begin{equation}
 u=s a_*,\qquad X=s^2Y_*,\qquad a_*|_S=\kappa.
 \label{var:zero-lapse-boundary}
\end{equation}
No assertion of a Lorentzian extension across the zero-lapse
boundary is made at this stage.
\end{proposition}

\begin{proof}
\emph{Smoothness and boundary moments.}
Proposition~\ref{var:area-elimination} removes the interior
surface measure from \eqref{var:hessian-measure}.  Together
with the differentiated KID equation, the resulting equations
express all second derivatives of $u,X$ as smooth linear
expressions in their first jets.  Ordinary elliptic
regularity makes them smooth in the open exterior.  The full
first-jet system is closed; along the normals of a smooth
collar it is a linear ODE with smooth coefficients.  Starting
on an interior collar face, solve that ODE to $S$.
Smooth dependence on the tangential parameters gives smooth
extension of all fields and their jets.  This argument also
shows that a zero full first jet at one interior point would
force the solution to vanish throughout the connected open
exterior.

There can still, at this point, be constraint moment measures
supported on $S$.  Integrate the smooth interior fields by
parts in \eqref{var:separation}.  Take $K'=0$ and both form
variations zero, and choose a metric test with value and
first jet zero at $S$.  The trace of its normal second
derivatives in the tangential slots is arbitrary, so its
scalar-curvature variation at $S$ is arbitrary.  Its area
and expansion variations, all integrated interior boundary
terms, momentum variation and ray-transport terms at $S$
vanish.  Equation~\eqref{var:separation} therefore kills
the scalar boundary moment.  Positivity and domination of
the vector moment by the scalar moment kill the vector
boundary moment as well.  Henceforth the moments in that
identity are just the smooth functions $u,X$ times volume.

\emph{The end constants and positive lapse.}
The closed first-jet system gives, for $Y=(u,X)$ in end
coordinates,
\begin{equation}
 |D^2Y|\le C r^{-1-q}|DY|+C r^{-2-q}|Y|.
 \label{var:end-jet-estimate}
\end{equation}
Indeed curvature, $\nabla K$, $K^2$, the constraints and
the electric stress multiply values, whereas $K$ and
metric first derivatives multiply first derivatives.
These orders require only the stated $O_2/O_1$ tails.
Radial integration and Gronwall first give $Y=O(r)$ and
$DY=O(1)$: write $|Y(r)|$ as its initial value plus the
radial integral of $|DY|$ and use the integrability of
$r^{-1-q}$.  Equation~\eqref{var:end-jet-estimate} now
gives a limit of $DY$ on each ray with error $O(r^{-q})$.
The limits agree on all rays, since integration of $D^2Y$
along a path of length $O(r)$ on a large sphere gives an
$O(r^{-q})$ difference.  Denote the common linear part by
$\mathcal A$.

The scalar affine part cannot be nonzero: $u\ge0$ on all
large spheres.  Once that part is zero, $|X|\le u$ rules
out every growing vector affine part as well.  Thus
$\mathcal A=0$.  Since $n\ge4$ and $q>(n-2)/2$, we have
$q>1$.  Integration gives bounded $Y$, a common constant
limit $Y_\infty$, and initially $Y-Y_\infty=O(r^{1-q})$.
Reinserting boundedness in \eqref{var:end-jet-estimate}
improves the Hessian to $O(r^{-2-q})$, and two integrations
give $Y-Y_\infty=O_2(r^{-q})$.  The stated weaker exponent
$q_0$ is therefore available.

Apply \eqref{var:separation} to the end prototypes and
integrate by parts on large truncations.  The interior
adjoint equations and the compact boundary terms cancel.
Only the constant scalar multiplying the linear scalar
ADM flux, and twice the constant vector multiplying the
linear momentum flux, remain.  All other boundary terms
vanish by \eqref{var:end-jet-estimate}, the prototype
orders and the background decay.  The surviving identity
is
\begin{equation}
 2k\omega\bigl(b^0\delta\mathcal E_{\mathrm{ADM}}
                     +b^i\delta(\mathbf P_{\mathrm{ADM}})_i\bigr)
 =2k\omega\bigl(u_\infty\delta\mathcal E_{\mathrm{ADM}}
                     +X_\infty^i\delta(\mathbf P_{\mathrm{ADM}})_i\bigr).
 \label{var:prototype-normalization}
\end{equation}
The prototype fluxes span all entries, proving
\eqref{var:lapse-shift-end}.  If $u$ vanished at an
interior point, smooth nonnegativity would give
$\dd u=0$ there.  Domination would give $X=0$ and
$\nabla X=0$ there as well, since $|X|\le u=O(d^2)$.
The full-jet uniqueness noted above would contradict
$u_\infty=b^0>0$.  Thus $u>0$ in the interior.

\emph{The stationary Einstein equation.}
The stationary metric in \eqref{var:stationary-product}
has second fundamental form
$b=-\operatorname{sym}\nabla X/u=K$ on each slice.
We recall the metric-adjoint computation, now with no
interior area source.  Put
$\mathcal Q(A,B)=A:B-(\operatorname{tr}A)(\operatorname{tr}B)$.
Integration of the momentum pairing gives
\begin{equation}
 u\bigl(R_g-\mathcal Q(K,K)+2\mathcal Q(K,b)\bigr)
 =u\bigl(R_g+\mathcal Q(b,b)-\mathcal Q(K-b,K-b)\bigr).
 \label{var:stationary-action}
\end{equation}
At $b=K$ its metric first variation equals that of the
stationary scalar action, at fixed lapse and shift, up to
the integrated divergence.  Its spatial coefficient is
$-u\operatorname{Ein}_{ij}$.
Equation~\eqref{var:electric-metric-adjoint} supplies
the electromagnetic coefficient.  The variation of the
constraint volume factor vanishes by complementarity,
and isometric transport of the active unit-ball rays
adds $-(J_m)_{(i}X_{j)}$.  Thus
\begin{equation}
 u(\operatorname{Ein}-\mathsf S)_{ij}
                         =-(J_m)_{(i}X_{j)}.
 \label{var:spatial-einstein-adjoint}
\end{equation}
Gauss--Codazzi gives the remaining normal and mixed
components as $\mu_m,J_m$.  Where $\mu_m=0$, DEC gives
$J_m=0$.  Where $\mu_m>0$, equality in
\begin{equation*}
 -J_m(X)\le |J_m||X|\le\mu_m u
 \end{equation*}
forces $|X|=u$ and $J_m=-\mu_m X^\flat/u$.
Consequently all components combine to give
\eqref{var:stationary-einstein}, and the extra matter
can occur only at $W=0$.

\emph{Maxwell tests and a global potential.}
Test \eqref{var:separation} with the metric and second
tensor fixed and a compactly supported closed variation
of $\alpha$.  The resulting identity, after dividing
by the nonzero factor $4c_n^2$, is
\begin{equation}
 \int_\Omega uE^\flat\wedge\alpha'=0
 \quad\hbox{for every compactly supported closed }
                         (n-1)\hbox{-form }\alpha'.
 \label{var:electric-form-test}
\end{equation}
Exact tests imply $\dd(uE^\flat)=0$.
All closed tests, not only exact ones, are available.
Compact-support de Rham duality on the oriented open
manifold $\operatorname{int}\Omega$ therefore makes
$uE^\flat$ exact; see \cite{BottTu1982}.  This step does
not presume simple connectivity.  Define $p$ by
$\dd p=(k-1)uE^\flat$.  Since $u$ is bounded at infinity
and $E=O_1(r^{-k})$, radial integration and integration
along large spheres give one common limit of $p$ at
infinity and $p-p_\infty=O_2(r^{2-n})$.
Choose $p_\infty=0$.  Integration in local collars
extends $p$ smoothly to $S$.

Independent closed compact two-form tests of $D_m$ give
\begin{equation}
 \int_\Omega D_m'\wedge i_X\alpha=0,
 \qquad\hbox{and hence}\qquad \dd(i_X\alpha)=0.
 \label{var:magnetic-form-test}
\end{equation}
The degree of $i_X\alpha$ is $n-2$, so this pairing
and its differential have the required degrees in
every dimension.  With orientation
$\vol_{\mathbf g}=u\dd z\wedge\vol_g$, a direct
orthonormal-frame computation gives
\begin{equation}
 *_\mathbf g\mathbf F=-c_n
          \bigl(\alpha+\dd z\wedge i_X\alpha\bigr).
 \label{var:star-field}
\end{equation}
Equations~\eqref{var:electric-form-test},
\eqref{var:magnetic-form-test} and $\dd\alpha=0$
therefore give both Maxwell equations.  Contraction
of $\mathbf F$ with $N=(\partial_z-X)/u$ and pullback
to a slice yields $c_nE^\flat$; its spatial pullback
is zero.  This verifies all conventions in the
claimed original-data realization.

\emph{Boundary laws.}
The outward integration normal of $\Omega$ at $S$
is $-\nu$.  The boundary coefficient of a free
second-tensor variation in \eqref{var:separation} is
\begin{equation}
 -2\operatorname{sym}(X^\flat\otimes\nu^\flat)
       +2X_\nu g
       -\frac{\dd\zeta}{\dd A_g}
                         (g-\nu^\flat\otimes\nu^\flat).
 \label{var:boundary-pairing-k}
\end{equation}
Interpreting this first as an equality of measures,
its tangential and mixed entries show
$X_{\rm tan}=0$ and $\dd\zeta=2X_\nu\dd A_g$.
Now use compact scalar tests with
$(h,K')=(2s g,sK)$ and the forms fixed.  Their
constraint differential terms are
$-2ku\Delta s+2kK(X,\nabla s)$; the other terms
cancel by the interior adjoint equation.  Since
$\theta_+=0$, their expansion derivative is
$k\partial_\nu s$, while their area derivative is
$k\int_Ss$.  Thus
\begin{equation}
 -kc\int_Ss\,\dd A_g
 =2k\int_S\left[u\partial_\nu s
       -(\partial_\nu u+K(X,\nu))s\right]\dd A_g
             -k\int_S\partial_\nu s\,\dd\zeta.
 \label{var:boundary-conformal-pairing}
\end{equation}
The boundary value and normal derivative of $s$ are
independent.  The derivative coefficient gives
$u=X_\nu$, and the value coefficient gives
$\partial_\nu u+K(X,\nu)=c/2$.

For the last boundary identity, take Lie variations
of all the data along a vector field extending
arbitrary $s\nu$ on $S$.  Such tests are legitimate
field variations on the fixed manifold: the jets
may be defined by a smooth extension across $S$;
no flow preserving the boundary is asserted.
The active interior constraints have zero first
variation by spatial transport at a minimum of
their nonnegative contractions.  The induced
rotation of the transported unit ball contributes
zero for the same reason.  The area and expansion
derivatives at $S$ are, respectively, $\int_SHs$
and $L_+s$.  Hence
\begin{equation*}
 c\int_SHs\,\dd A_g=2\int_SuL_+s\,\dd A_g,
 \end{equation*}
which proves $cH=2L_+^*u$.
One-sided outer area minimality gives $H\ge0$ by
nonnegative outward tests.  The strong minimum
principle for this smooth elliptic equation and
$u\ge0$ on connected $S$ gives the boundary
dichotomy in the proposition.

The boundary laws give $W=0$ on $S$.  If $u>0$
there, differentiate $W$ and use
\eqref{var:first-kid}; tangential derivatives
vanish and
\begin{equation*}
 \partial_\nu W
   =2u\bigl(\partial_\nu u+uK(\nu,\nu)\bigr)
   =2\kappa u.
\end{equation*}
This is \eqref{var:positive-lapse-boundary} and
implies positivity in a deleted exterior collar.
If $u=0$ on $S$, then $X=0$ there and
$\partial_\nu u=\kappa$.  All tangential derivatives
of $X$ vanish at $S$; the normal-normal and mixed
parts of the KID equation then give
$\nabla_\nu X=0$ there.  Smooth division by the
boundary defining function gives
\eqref{var:zero-lapse-boundary}.  In particular
$W=\kappa^2s^2+O(s^3)>0$ in a deleted collar.
Finally \eqref{var:lapse-shift-end} and
$(b^0)^2-|b|^2=1$ imply $W\to1$.
\end{proof}

On the positive set of $W$ it will be useful to introduce
\begin{equation}
\begin{gathered}
 \lambda=\sqrt W,\qquad
 h_s=g+W^{-1}X^\flat\otimes X^\flat,\\
 C=h_s^{-1}=g^{-1}-u^{-2}X\otimes X,\qquad
 A_s=W^{-1}X^\flat.
\end{gathered}
 \label{var:orbit-metric}
\end{equation}
The stationary metric is then
$\mathbf g=-\lambda^2(\dd z-A_s)^2+h_s$.
These definitions are made only where $W>0$; the next
section proves global staticity and rules out interior
zeros, rather than imposing either conclusion here.
